\section{Conclusion}
\label{section:conclusion}

This paper set out to do two things at once: survey the small-language-model
landscape as it stands in 2026, and argue a position about where it is going. On
the first, we cataloged 125 documented SLM variants under a four-dimensional taxonomy,
tracked the 2025--2026 shift toward hybrid state-space architectures, unified
reasoning models, and mobile-first multimodal families, reviewed the compression
techniques specific to the sub-10B regime---where quantization consistently
outperforms pruning \cite{kurz2025pruning_vs_quant}---and consolidated publicly
reported performance for a 17-model cohort across 15 tasks, six hardware tiers,
and three precisions, with every figure traced to a dated source rather than an
original run.

On the second: the ``less is more'' thesis is no longer only about efficiency. The
consumer of language intelligence that is growing fastest is the AI agent, and the
agent does not want a generalist. It issues many narrow, repetitive,
format-constrained calls, and for those calls a small specialized model is
sufficiently capable, 10--30$\times$ cheaper to serve, cheaper to adapt, and
easier to hold to a code interface \cite{belcak2025slm_agentic}. The consolidated
evidence in this paper supports each premise of that argument, and two
production-scale deployments---Apple's on-device/server split and the
routing-based systems surveyed in Section~\ref{sec:agentic}---already embody it.
The barriers that remain are those of sunk infrastructure, benchmark
mis-measurement, and attention, not of capability.

The work ahead is concrete. Reasoning still shows a performance cliff below
roughly 7B parameters that inference-time compute only partly closes;
hallucination is worse at small scale; compression can amplify bias; hardware and
model architecture are still designed separately; and the consolidated evidence
base needs controlled verification of the handful of claims that practitioners
most depend on (Section~\ref{sec:challenges}). None of these is a reason to wait.
The future of applied NLP is not one frontier model answering every query through
a cloud API. It is a fabric of small, specialized, mostly on-device models doing
the ordinary work, calling a large model only when the task genuinely demands it.
In that architecture, less is not a limitation. It is the design.
